% C055 — 原有结论保留，并把适用条件及切线推广写完整。
\begin{statementbox}
设抛物线 $y^2=2px\ (p>0)$，焦点 $F(p/2,0)$。经过 $F$ 的直线与曲线交于两个不同的有限点 $A(x_1,y_1),B(x_2,y_2)$，则 $AB$ 为焦点弦。
\begin{center}\begin{tikzpicture}[font=\small,>=stealth]
\draw[rounded corners=3pt,gray!45] (0,0) rectangle (14.2,10);
\node at (7.1,9.45) {先看懂对象：一条弦穿过焦点，两端点分居横轴两侧};
\begin{scope}[shift={(1.65,3.15)},scale=1.05]
\draw[->,gray] (-.65,0)--(5.3,0) node[above] {$x$};
\draw[->,gray] (0,-2.05)--(0,4.95) node[left] {$y$};
\draw[teal!75!black,thick,samples=140,domain=-1.85:4.55,variable=\t]
 plot ({\t*\t/4},{\t});
\draw[blue!65!black,dashed] (0,4)--(4,4)--(4,0);
\draw[orange!85!black,dashed] (0,-1)--(.25,-1)--(.25,0);
\draw[blue!75!black,very thick] (4,4)--(1,0);
\draw[orange!85!black,very thick] (1,0)--(.25,-1);
\fill (4,4) circle(2pt) node[above left=7pt] {$A$};
\fill (.25,-1) circle(2pt) node[right=14pt] {$B$};
\fill (1,0) circle(2pt) node[below right=7pt] {$F$};
\node[above left=7pt] at (0,0) {$O$};
\node[blue!75!black,left=8pt] at (0,4) {$y_1$};
\node[orange!85!black,left=8pt] at (0,-1) {$y_2$};
\fill (4,0) circle(1.3pt) node[below=8pt,blue!75!black] {$x_1$};
\fill (.25,0) circle(1.3pt) node[above right=9pt,orange!85!black] {$x_2$};
\end{scope}
\node[anchor=west,align=left] at (8.6,6.3)
 {$\textcolor{blue!75!black}{A(x_1,y_1),\quad y_1>0}$\\[13pt]
 $\textcolor{orange!85!black}{B(x_2,y_2),\quad y_2<0}$\\[16pt]
 $A,F,B$ 三点共线\\[13pt]
 $F\!\left(\dfrac p2,0\right)$ 在弦内\\[16pt]
 两端不必上下对称};
\node at (7.1,1.1) {横坐标同号，积为正；纵坐标异号，积为负。};
\node at (7.1,.45) {图取 $p=2$ 的一条斜弦；坐标积定值适用于任意有限焦点弦。};
\end{tikzpicture}
\end{center}

\medskip
\textbf{\textcolor{CStatement}{两条必记定值}}\par\smallskip
\[
 \boxed{\ x_1x_2=\frac{p^2}{4}\ },\qquad
 \boxed{\ y_1y_2=-p^2\ }.
\]
两端点可以不对称；只要是有限焦点弦，上面两个坐标积就不随方向改变。

\end{statementbox}
\clearpage
\begin{statementbox}
\textbf{\textcolor{CStatement}{新推论：顶点看焦点弦，两侧张角的正切积为 $4$}}\par\smallskip
\begin{center}% Exact model: p=2, A=(4,4), B=(1/4,-1), F=(1,0).
\begin{tikzpicture}[font=\small,>=stealth]
\draw[rounded corners=3pt,gray!45] (0,0) rectangle (14.2,8.8);
\node at (7.1,8.25) {从顶点看焦点弦：两个锐角的正切积为正};
\begin{scope}[shift={(1.5,3.0)},scale=.9]
\fill[blue!6] (0,0)--(4,0)--(4,4)--cycle;
\fill[orange!9] (0,0)--(.25,0)--(.25,-1)--cycle;
\draw[->,gray] (-.6,0)--(5.8,0) node[above] {$x$};
\draw[->,gray] (0,-1.8)--(0,4.8) node[left] {$y$};
\draw[teal!75!black,thick,samples=120,domain=-1.65:4.6,variable=\t]
 plot ({\t*\t/4},\t);
\draw[gray!70,dashed,thick] (4,4)--(.25,-1);
\draw[blue!75!black,very thick] (0,0)--(4,4);
\draw[orange!85!black,very thick] (0,0)--(.25,-1);
\draw[blue!75!black,dashed] (4,4)--(4,0);
\draw[violet,thick] (.85,0) arc[start angle=0,end angle=45,radius=.85];
\node[violet] at (1.7,.4) {$\alpha$};
\draw[violet,thick] (1.5,0) arc[start angle=0,end angle=-75.9637565,radius=1.5];
\node[violet] at (1.48,-1.13) {$\beta$};
\fill (0,0) circle(2pt) node[above left=7pt] {$O$};
\fill (1,0) circle(2pt) node[below=8pt] {$F$};
\fill (4,4) circle(2pt) node[above left=6pt] {$A(4,4)$};
\fill (.25,-1) circle(2pt);
\node[anchor=east] at (-.35,-1.45) {$B$};
\draw[gray,thin] (-.3,-1.25)--(.18,-1.03);
\end{scope}
\node[anchor=west,align=left] at (8.0,4.5)
 {$y^2=4x,\quad F(1,0)$\\[12pt]
 $B(\frac14,-1)$\\[12pt]
 $\alpha=\angle AOF,\ \beta=\angle FOB$\\[14pt]
 $\tan\alpha=1,\quad\tan\beta=4$\\[14pt]
 $\tan\alpha\tan\beta=4$\\[14pt]
 $k_{OA}=1,\quad k_{OB}=-4$};
\node at (7.1,.65) {蓝、橙实线从顶点出发；灰色虚线是过焦点的弦 $AB$。};
\end{tikzpicture}
\end{center}
设顶点为 $O(0,0)$，$A$ 在 $x$ 轴上方，$B$ 在下方，令
$\alpha=\angle AOF$、$\beta=\angle FOB$，两角均为锐角，则
\[
 \boxed{\tan\alpha\,\tan\beta=4}.
\]
等价地，两条顶点连线 $OA,OB$ 的\textbf{带符号斜率积为 $-4$}：
$k_{OA}k_{OB}=-4$。这里的 $4$ 是\textbf{两个角的正切积}，不是 $\tan\angle AOB$。
\end{statementbox}
\clearpage
\begin{statementbox}
\textbf{\textcolor{CStatement}{更一般：过对称轴上定点的弦，顶点两侧正切积}}\par\smallskip
设 $P(a,0)$ 位于抛物线对称轴的\textbf{正半轴}，其中 $a>0$。
过 $P$ 的直线与 $y^2=2px$ 交于 $A(x_1,y_1)$、$B(x_2,y_2)$，
按 $y_1>0>y_2$ 编号；顶点为 $O(0,0)$，令
$\alpha=\angle AOP$、$\beta=\angle POB$（均为锐角），则
\[
 \boxed{\tan\alpha\,\tan\beta=\frac{2p}{a}}.
\]
同时有两个坐标积定值及带符号斜率积：
\[
 \boxed{y_1y_2=-2pa},\qquad
 \boxed{x_1x_2=a^2},\qquad
 \boxed{k_{OA}k_{OB}=-\frac{2p}{a}}.
\]
\begin{center}% Exact model: p=2, a=2, x=y/2+2, A=(4,4), B=(1,-2).
\begin{tikzpicture}[font=\small,>=stealth]
\draw[rounded corners=3pt,gray!45] (0,0) rectangle (14.2,9.6);
\node at (7.1,9.05) {过轴上正半轴定点的弦：从两个投影三角形读取正切};
\begin{scope}[shift={(1.45,3.9)},scale=.85]
\fill[blue!6] (0,0)--(4,0)--(4,4)--cycle;
\fill[orange!9] (0,0)--(1,0)--(1,-2)--cycle;
\draw[->,gray] (-.65,0)--(5.8,0) node[above] {$x$};
\draw[->,gray] (0,-2.8)--(0,4.95) node[left] {$y$};
\draw[teal!75!black,thick,samples=120,domain=-2.5:4.65,variable=\t]
 plot ({\t*\t/4},\t);
\draw[gray!75,dashed,thick] (4,4)--(1,-2);
\draw[blue!75!black,very thick] (0,0)--(4,4);
\draw[orange!85!black,very thick] (0,0)--(1,-2);
\draw[blue!75!black,dashed] (4,4)--(4,0) node[midway,right=7pt] {$y_1$};
\draw[orange!85!black,dashed] (1,-2)--(1,0);
\node[orange!85!black,anchor=west] at (1.4,-1.8) {$-y_2$};
\draw[gray,thin] (1.35,-1.65)--(1,-1.3);
\draw[gray] (3.8,0)--(3.8,.2)--(4,.2);
\draw[gray] (.8,0)--(.8,-.2)--(1,-.2);
\draw[violet,thick] (.8,0) arc[start angle=0,end angle=45,radius=.8];
\node[violet] at (1.8,1.1) {$\alpha$};
\draw[violet,thick] (1.45,0) arc[start angle=0,end angle=-63.4349488,radius=1.45];
\node[violet] at (2.2,-.8) {$\beta$};
\fill (0,0) circle(2pt) node[above left=7pt] {$O$};
\fill (2,0) circle(2pt) node[above right=8pt] {$P(a,0)$};
\fill (4,4) circle(2pt) node[above left=6pt] {$A$};
\fill (1,-2) circle(2pt) node[below left=7pt] {$B$};
\node[blue!75!black,below=8pt] at (4,0) {$x_1$};
\node[orange!85!black,above=8pt] at (1,0) {$x_2$};
\end{scope}
\node[anchor=west,align=left] at (8.2,5.8)
 {图取 $p=2,\ a=2$\\[12pt]
 $A(4,4),\quad B(1,-2)$\\[14pt]
 $\tan\alpha=\dfrac{y_1}{x_1}=1$\\[15pt]
 $\tan\beta=\dfrac{-y_2}{x_2}=2$\\[15pt]
 $\tan\alpha\tan\beta=\dfrac{2p}{a}=2$};
\node at (7.1,1.0) {$x_1,x_2>0,\ y_1>0>y_2$：两角为锐角，两侧正切都取正值。};
\node at (7.1,.4) {弦写成 $x=my+a$；本图 $m=1/2$，竖直弦对应 $m=0$。};
\end{tikzpicture}
\end{center}
\textbf{焦点弦是特例：}当 $P=F$ 时 $a=\frac p2$，立即得到
$\tan\alpha\tan\beta=\frac{2p}{p/2}=\boxed{4}$。
\textbf{条件提醒：}$a>0$ 是本条“两个正角正切积”公式的必要设定；
$a=0$ 时可能有端点为顶点，正切比无法定义；$a<0$ 时两端点不再一上一下，不能原样套用 $2p/a$。
\end{statementbox}
\clearpage
\begin{statementbox}

\medskip
\textbf{\textcolor{CStatement}{推广：焦点弦两端点切线互相垂直}}\par\smallskip
\begin{center}\begin{tikzpicture}[font=\small,>=stealth]
\draw[rounded corners=3pt,gray!45] (0,0) rectangle (14.2,12.5);

\node at (7.1,11.95) {切线推广：两切线成直角，交点恰好落在准线上};
\begin{scope}[shift={(3.1,4.2)},scale=1.05]
\draw[->,gray] (-1.65,0)--(4.9,0) node[above] {$x$};
\draw[->,gray] (0,-2.2)--(0,5.15) node[left] {$y$};
\draw[violet,dashed] (-1,-2.2)--(-1,5.05);
\node[violet,rotate=90] at (-1.5,3.8) {准线 $x=-1$};
\draw[teal!75!black,thick,samples=120,domain=-2.2:4.5,variable=\t] plot ({\t*\t/4},{\t});
\draw[gray!75,thick] (4,4)--(.25,-1);
\draw[blue!75!black,very thick] (-1,1.5)--(4,4);
\draw[orange!85!black,very thick] (-1,1.5)--(.25,-1);
\draw[black] (-.68695,1.65652)--(-.53043,1.34348)--(-.84348,1.18695);
\fill (-1,1.5) circle(2pt) node[left=8pt] {$T$};
\fill (4,4) circle(2pt) node[above left=7pt] {$A$};
\fill (.25,-1) circle(2pt) node[below left=7pt] {$B$};
\fill (1,0) circle(2pt) node[below right=6pt] {$F$};
\end{scope}
\node[anchor=west,align=left] at (9.1,7.15)
 {$A(4,4)$\\[12pt]$B(1/4,-1)$\\[12pt]$T(-1,3/2)$\\[16pt]
 蓝色切线：$k_A=1/2$\\[12pt]橙色切线：$k_B=-2$\\[14pt]$k_Ak_B=-1$};
\node at (7.1,1.4) {灰线是焦点弦 $AB$；蓝、橙两线是端点处的切线。};
\node at (7.1,.55) {$y^2=4x,\quad p=2,\quad y_Ay_B=-4\quad\Longrightarrow\quad k_Ak_B=-1$};
\end{tikzpicture}
\end{center}
分别在 $A,B$ 处作抛物线的切线，切线斜率为
\[
 k_A=\frac{p}{y_1},\qquad k_B=\frac{p}{y_2},\qquad
 \boxed{k_Ak_B=-1}.
\]
两条切线交于 $T$，其交点\textbf{必在准线} $x=-\frac p2$ 上：
\[
 \boxed{\ T\!\left(-\frac p2,\frac{y_1+y_2}{2}\right)\ }.
\]

\medskip
\textbf{\textcolor{CStatement}{范围与变形}}\par\smallskip
\begin{itemize}[leftmargin=*,itemsep=0.2em]
 \item 直线 $y=0$ 虽过焦点，却只与抛物线交于顶点 $O$，\textbf{不构成两点焦点弦}；其余过焦点的直线均有两个不同交点。
 \item 对左开口 $y^2=-2px$（$p>0$），两个坐标积仍依次为 $\frac{p^2}{4}$、$-p^2$。
 \item 对上下开口 $x^2=\pm2py$（$p>0$），\textbf{交换 $x,y$}：$x_1x_2=-p^2$、$y_1y_2=\frac{p^2}{4}$。
\end{itemize}
\end{statementbox}
